\section{中心化坐标、归一化恢复与损失分解}
\label{sec:inner-quotient}
对任意正分布 $\mu$，定义 $C_\mu h=h-\E_\mu h$。这是到均值为零子空间的线性投影：
\begin{equation}
C_\mu^2=C_\mu,\qquad C_\mu\mathbf1=0,\qquad
\ker C_\mu=\operatorname{span}\{\mathbf1\}.
\end{equation}
后一个等式由 $C_\mu h=0\Longleftrightarrow h_a=\E_\mu h$ 对全部 $a$ 成立得到。商空间把相差常数的 log-ratio 视为同一个坐标，但概率归一化仍唯一确定该常数。

\begin{theorem}[归一化分布与中心化误差的一一对应]
\label{thm:quotient-bijection}
给定 $p^{\mathrm G}>0$ 和 $q>0$，令
\begin{equation}
e=C_q\log\frac{p}{p^{\mathrm G}},\qquad
A(e)=\log\E_{p^{\mathrm G}}e^{e}.
\end{equation}
映射 $p\mapsto e$ 是 $\Delta_V^\circ$ 与 $\{e:\E_q e=0\}$ 的双射，其逆映射为
\begin{equation}
p_a(e)=p_a^{\mathrm G}\exp(e_a-A(e)).
\label{eq:quotient-inverse}
\end{equation}
\end{theorem}
\begin{proof}
令 $h_a=\log(p_a/p_a^{\mathrm G})$，$c=\E_qh$，于是 $h=e+c$。利用 $p$ 归一化，
\begin{align}
1=\sum_a p_a&=e^c\sum_a p_a^{\mathrm G}e^{e_a},\\
c&=-\log\sum_a p_a^{\mathrm G}e^{e_a}=-A(e).
\end{align}
这证明任意 $p$ 满足式~\eqref{eq:quotient-inverse}。反过来，任意 $\E_q e=0$ 的向量经该式构造的 $p$ 为正且和为一，并满足
\begin{equation}
C_q\log(p/p^{\mathrm G})=C_q(e-A(e)\mathbf1)=e.
\end{equation}
两方向复合均为恒等映射，故双射成立。
\end{proof}

\begin{theorem}[四种目标在同一坐标中的精确表达]
\label{thm:quotient-losses}
设 $G=\eta C_qf$，$Z_f=\E_qe^{\eta f}$，并记
\begin{equation}
\Delta=\log Z_f-\eta\E_q f
=\log\E_q e^G=\KL(q\Vert p^{\mathrm G}).
\end{equation}
在上一双射坐标中，三个残差及四种目标为
\begin{align}
R^{\mathrm{Part}}&=e-A(e),&
R^{\mathrm{BPO}}&=e,&
R^{\mathrm{Mean}}&=e-A(e)-\Delta,\\
\mathcal L_{\mathrm{Part}}&=\tfrac12\E_qe^2+\tfrac12A(e)^2,\\
\eta\mathcal L_{\mathrm{BPO}}&=\tfrac12\E_qe^2,\\
\mathcal L_{\mathrm{Mean}}&=\tfrac12\E_qe^2+\tfrac12(A(e)+\Delta)^2,\\
\mathcal L_{\mathrm{MDPO}}&=\frac1\eta\KL(p(e)\Vert p^{\mathrm G})-\frac1\eta\log Z_f.
\label{eq:quotient-four-losses}
\end{align}
因此 $\log Z_f$ 在 BPO 中消去后，仍由 $p(e)$ 的归一化恢复。PMD-Part 与 BPO 有相同的零误差分布，却具有不同的目标函数。
\end{theorem}
\begin{proof}
Gibbs 方程给出 $\log(p^{\mathrm G}/q)=G-\Delta$。把它和双射式相加，得到
\begin{equation}
U=\log(p/q)=G-\Delta+e-A(e).
\end{equation}
于是 $R^{\mathrm{Part}}=U-\eta f+\log Z_f=e-A(e)$，$R^{\mathrm{Mean}}=U-G=e-A(e)-\Delta$。中心化完整残差得
\begin{equation}
C_qR^{\mathrm{Part}}=C_q(e-A(e))=e=R^{\mathrm{BPO}}.
\end{equation}
平方展开并使用 $\E_qe=0$，有
\begin{align}
\E_q(e-A)^2&=\E_qe^2-2A\E_qe+A^2=\E_qe^2+A^2,\\
\E_q(e-A-\Delta)^2&=\E_qe^2+(A+\Delta)^2.
\end{align}
MDPO 表达由 Gibbs 变分恒等式给出。最后按 $q$ 对 $U$ 求均值，
\begin{equation}
\E_qU=-\Delta-A(e),\qquad \KL(q\Vert p)=\Delta+A(e)\ge0.
\end{equation}
这也证明 $A(e)$ 本身可以为负，其下界是 $-\Delta$；它的定义使用 $p^{\mathrm G}$ 求期望，$e$ 却按 $q$ 中心化，二者不可混用。
\end{proof}

直接使用平方和中心化的正交分解，还可写成
\begin{align}
\mathcal L_{\mathrm{Part}}&=\eta\mathcal L_{\mathrm{BPO}}
+\tfrac12\bigl[\Delta-\KL(q\Vert p)\bigr]^2,\\
\mathcal L_{\mathrm{Mean}}&=\eta\mathcal L_{\mathrm{BPO}}
+\tfrac12\KL(q\Vert p)^2.
\label{eq:quotient-penalties}
\end{align}
第一行保留完整残差的均值误差；第二行把 student 对 anchor 的 KL 加入额外惩罚。两者分别解释相同 Gibbs 最优点下的轨迹差异，以及 Mean 改变充分优化后分布的原因。

\begin{proposition}[成对回归完全消去两个归一化常数]
\label{prop:quotient-pairwise}
令 $h_a=U_a-\eta f_a$，两个 action 独立来自 $q$。则
\begin{equation}
\mathcal L_{\mathrm{BPO}}=
\frac1{4\eta}\E_{a,b\sim q}(h_a-h_b)^2.
\end{equation}
若 $p=\operatorname{softmax}(z)$，差值为
\begin{equation}
h_a-h_b=z_a-z_b-\log q_a+\log q_b-\eta(f_a-f_b).
\end{equation}
\end{proposition}
\begin{proof}
独立性给出 $\E h_ah_b=(\E_qh)^2$，故
\begin{equation}
\E(h_a-h_b)^2=2\E_qh^2-2(\E_qh)^2=2\Var_qh.
\end{equation}
而 $\Var_qh=\E_q(R^{\mathrm{BPO}})^2$。再将 $\log p_a=z_a-\log\sum_v e^{z_v}$ 代入两个 $U$，两份 student 归一化常数也相消，得到第二个式子。概率归一化由 softmax 保持，Gibbs 配分函数无需作为数值 target 输入。
\end{proof}
